% Compiles on its own: latexmk "60 - evaluation" (in this directory, see README.md).
% The subfiles package loads the preamble of the main file named here.
\documentclass[../main-german]{subfiles}
%\documentclass[../main-basic-listings-german]{subfiles}
%\documentclass[../main-german-thu]{subfiles}
%\documentclass[../main-german-university-of-hamburg]{subfiles}
\begin{document}
%%% ===============================================================================
\chapter{Evaluation}\label{sec:evaluation}
%%% ===============================================================================
\progress{0}

\emph{Purpose and scope}.
The evaluation chapter should be the most formal and rigorously structured chapter of your thesis as the validity of your evaluation argument depends on it.

% ---------------------------------------------------------------------------------
\section{Objective (SM2)}\label{sec:evaluation:objective}
% ---------------------------------------------------------------------------------

State what you want to evaluate and what you want to measure.

% ---------------------------------------------------------------------------------
\section{Setup (SM3)}\label{sec:evaluation:setup}
% ---------------------------------------------------------------------------------

State which data, participants, tools, etc. you chose and why.
State how you measure outcomes and how you compare them to baselines, reference groups, etc.

% ---------------------------------------------------------------------------------
\section{Execution (SM4)}\label{sec:evaluation:execution}
% ---------------------------------------------------------------------------------

Provide all details on the execution that are necessary to allows another person to reproduce your results at a later point.

% ---------------------------------------------------------------------------------
\section{Results (SM5)}\label{sec:evaluation:results}
% ---------------------------------------------------------------------------------

You only report the measurements.
You must present and reference them (e.g., ``\zcref{fig:my_label2} shows the runtime'') and describe the results in text.
If you use diagrams, tables, or charts, you cannot expect the reader to know what to you expect them to see in a diagram, table or chart.
Describe to them how to read these, explain the meaning of particular elements, point out special observations.
Describe the results, but do not interpret them.
Make use of appendices if certain details take too much space.

\begin{figure}
  \centering
  %%%\includegraphics{/path/to/figure.pdf}
  \caption{Another scientific figure that has to be explained in the text}
  \label{fig:my_label2}
\end{figure}

% ---------------------------------------------------------------------------------
\section{Discussion (SM6)}\label{sec:evaluation:discussion}
% ---------------------------------------------------------------------------------

Interpret the results only after describing all of them.
The interpretation must be based on the written description in \zcref{sec:evaluation:results} only.
Then you can be sure that your readers can follow your interpretation and reach the same conclusions as you have.

% ---------------------------------------------------------------------------------
\section{Threats to Validity}\label{sec:evaluation:threats}
% ---------------------------------------------------------------------------------

Discuss what limits the validity of your results and how you mitigated it.
Typical questions are:

\begin{description}
  \item[Construct validity] Do your measurements capture what you intended to evaluate (\zcref{sec:evaluation:objective})?
  \item[Internal validity] Are the results caused by what you varied, or by other factors?
  \item[External validity] To which other settings, data, or participants do the results generalize?
  \item[Reliability] Would another person obtain the same results and conclusions following \zcref{sec:evaluation:execution}?
\end{description}
\end{document}
